\section*{Chapitre \ref{ch:g12:cells-and-electrolysis} --- Piles et électrolyse}

\begin{solution}{exo:g12:cells-and-electrolysis:1}
L'électrode de zinc est l'\omterm{def:g12:cells-and-electrolysis:anode}{anode}, \ce{Zn -> Zn^{2+} + 2e-}~; l'électrode
de cuivre est la \omterm{def:g12:cells-and-electrolysis:anode}{cathode}, \ce{Cu^{2+} + 2e- -> Cu}.
\end{solution}

\begin{solution}{exo:g12:cells-and-electrolysis:2}
% ledger: const:F
$Q = 0.20 \times 1800 = \qty{360}{C}$~; $n(e^-) = 360/96485 =
\qty{3.7e-3}{mol}$.
\end{solution}

\begin{solution}{exo:g12:cells-and-electrolysis:3}
Il ferme le circuit et maintient chaque solution électriquement
neutre, ses \omterm{def:g9:ions:ion}{ions} se déplaçant vers les compartiments. Sans lui, le
circuit est ouvert~: aucun courant ne circule, le voltmètre indique zéro
et une lampe reste éteinte.
\end{solution}

\begin{solution}{exo:g12:cells-and-electrolysis:4}
$2.5 \times 3600 = \qty{9000}{C}$.
\end{solution}

\begin{solution}{exo:g12:cells-and-electrolysis:5}
Non~: la réaction se fait dans le sens opposé au sens spontané. C'est le
générateur qui fournit l'énergie.
\end{solution}

\begin{solution}{exo:g12:cells-and-electrolysis:6}
\omterm{def:g12:cells-and-electrolysis:anode}{Anode} (cuivre, $-$)~: \ce{Cu -> Cu^{2+} + 2e-}. \omterm{def:g12:cells-and-electrolysis:anode}{Cathode} (argent, $+$)~:
\ce{Ag+ + e- -> Ag}. Réaction globale~: \ce{Cu + 2Ag+ -> Cu^{2+} + 2Ag}.
\end{solution}

\begin{solution}{exo:g12:cells-and-electrolysis:7}
Trois ans font $3 \times 365 \times 24 = \qty{26280}{h}$~;
$0.010 \times 26280 = \qty{263}{mA.h}$.
\end{solution}

\begin{solution}{exo:g12:cells-and-electrolysis:8}
% ledger: const:F, aw:Ag
$Q = 0.50 \times 2400 = \qty{1200}{C}$~; $n(e^-) = n(\ce{Ag}) =
1200/96485 = \qty{0.0124}{mol}$~; $m = 0.0124 \times 107.9 =
\qty{1.34}{g}$.
\end{solution}

\begin{solution}{exo:g12:cells-and-electrolysis:9}
Dans le fil, le courant va du cuivre ($+$) vers le zinc ($-$), en sens
opposé aux \omterm{def:g9:inside-the-atom:nucleus}{électrons}. Les \omterm{def:g9:ions:ion}{ions} \ce{K+} se déplacent vers le compartiment
du cuivre, qui perd des \omterm{def:g9:ions:ion}{ions} \ce{Cu^{2+}}~; les \omterm{def:g9:ions:ion}{ions} \ce{NO3-} vers le
compartiment du zinc, qui gagne des \omterm{def:g9:ions:ion}{ions} \ce{Zn^{2+}}.
\end{solution}

\begin{solution}{exo:g12:cells-and-electrolysis:10}
Dans les deux cas, l'\omterm{def:g12:cells-and-electrolysis:anode}{anode} est le siège de l'\omterm{def:g11:redox:oxidation}{oxydation}. Dans une \omterm{def:g12:cells-and-electrolysis:cell}{pile},
l'\omterm{def:g12:cells-and-electrolysis:anode}{anode} est la borne $-$, la réaction est spontanée et l'énergie
chimique devient énergie électrique~; dans une \omterm{def:g12:cells-and-electrolysis:electrolysis}{électrolyse}, l'\omterm{def:g12:cells-and-electrolysis:anode}{anode} est
reliée au $+$ du
générateur, la réaction est forcée et l'énergie électrique devient
énergie chimique.
\end{solution}

\begin{solution}{exo:g12:cells-and-electrolysis:11}
\qty{12}{mL}~: l'équation \ce{2H2O -> 2H2 + O2} donne deux \omterm{def:g7:atoms-and-molecules:molecule}{molécules} de
dihydrogène pour une de dioxygène, et des quantités égales de gaz
occupent des volumes égaux.
\end{solution}

\begin{solution}{exo:g12:cells-and-electrolysis:12}
% ledger: aw:Zn, const:F
$n(\ce{Zn}) = 5.0/65.4 = \qty{0.076}{mol}$~;
$n(\ce{Cu^{2+}}) = 0.100 \times 0.50 = \qty{0.050}{mol}$~; ils
réagissent \omterm{def:g10:the-mole:amount}{mole} à \omterm{def:g10:the-mole:amount}{mole}, donc les \omterm{def:g9:ions:ion}{ions} cuivre s'épuisent en premier.
$n(e^-) = 2 \times 0.050
= \qty{0.100}{mol}$, $Q = \qty{9.65e3}{C} = \qty{2.68}{A.h}$. À
\qty{50}{mA}~: $2.68/0.050 = \qty{54}{h}$.
\end{solution}

\begin{solution}{exo:g12:cells-and-electrolysis:13}
$Q = 2.0 \times 3600 = \qty{7200}{C}$~; $E = 7200 \times 1.5 =
\qty{1.08e4}{J}$, soit $1.5 \times 2.0 = \qty{3.0}{W.h}$.
\end{solution}

\begin{solution}{exo:g12:cells-and-electrolysis:14}
% ledger: const:F
$Q = \qty{7200}{C}$, $n(e^-) = \qty{0.0746}{mol}$. Deux \omterm{def:g9:inside-the-atom:nucleus}{électrons} par
\ce{H2}~: $n(\ce{H2}) = \qty{0.0373}{mol}$, $V = \qty{0.90}{L}$. Quatre
par
\ce{O2}~: $n(\ce{O2}) = \qty{0.0187}{mol}$, $V = \qty{0.45}{L}$.
\end{solution}

\begin{solution}{exo:g12:cells-and-electrolysis:15}
Le cuivre dissous à l'\omterm{def:g12:cells-and-electrolysis:anode}{anode} est égal au cuivre déposé, les mêmes
\omterm{def:g9:inside-the-atom:nucleus}{électrons} passant aux deux électrodes~: $2.84/3.00 = \qty{94.7}{\percent}$
de cuivre.
\end{solution}

\begin{solution}{pb:g12:cells-and-electrolysis:1}
% ledger: const:F, const:NA, aw:Li, dch:CH4
\textbf{1.} L'\omterm{def:g12:cells-and-electrolysis:anode}{anode}~: le lithium y est oxydé.

\textbf{2.} L'électrode de graphite, l'\omterm{def:g12:cells-and-electrolysis:anode}{anode}.

\textbf{3.} Les \omterm{def:g9:inside-the-atom:nucleus}{électrons} traversent le téléphone de l'électrode de
graphite vers l'électrode d'oxyde~; les \omterm{def:g9:ions:ion}{ions} lithium traversent
l'électrolyte dans le même sens, du graphite vers l'oxyde.

\textbf{4.} Sa réaction peut être forcée en sens inverse par un
chargeur.

\textbf{5.} \qty{4.000}{A.h}, soit $4.000 \times 3600 =
\qty{14400}{C}$.

\textbf{6.} $14400/96485 = \qty{0.1492}{mol}$.

\textbf{7.} $0.1492 \times \num{6.022e23} = \num{8.99e22}$ \omterm{def:g9:inside-the-atom:nucleus}{électrons}.

\textbf{8.} $4.000/0.25 = \qty{16}{h}$.

\textbf{9.} $0.20 \times 14400 = \qty{2880}{C}$.

\textbf{10.} $E = 14400 \times 3.7 = \qty{5.3e4}{J}$, soit environ
\qty{53}{kJ}.

\textbf{11.} $53280/3600 = \qty{14.8}{W.h}$.

\textbf{12.} La \omterm{def:g4:what-a-fire-needs:burning}{combustion} de \qty{1}{g} de méthane libère un peu plus
d'énergie (\qty{55.7}{kJ}) que n'en stocke une batterie de téléphone
pleine.

\textbf{13.} $1000/14.8 = 67.6$~: 67 charges complètes.

\textbf{14.} \ce{Li+ + e- -> Li}~: une \omterm{def:g11:redox:oxidation}{réduction}.

\textbf{15.} Le chargeur force la réaction dans le sens opposé au sens
spontané.

\textbf{16.} $14400/2.0 = \qty{7200}{s}$, soit \qty{2.0}{h}.

\textbf{17.} Un \omterm{def:g9:ions:ion}{ion} lithium par \omterm{def:g9:inside-the-atom:nucleus}{électron}~: \qty{0.1492}{mol}.

\textbf{18.} $0.1492 \times 6.9 = \qty{1.03}{g}$.

\textbf{19.} Environ \qty{1.03}{g} de lithium fait la navette entre les
électrodes au cours d'une charge complète.
\end{solution}
